\chapter{Completing access to stationary channel currents}
\label{ch:graph-access}

Conservation turns some inaccessible currents into measurable combinations of
other currents. It also has a definite limit: conservation of the aggregate
flow cannot determine how distinct physical channels divide that flow. This
chapter separates those two situations and gives exact sensor designs and
independently checkable uncertainty certificates.

The graph, its physical channels, and their assignment to pairs of states are
assumed known. The results do not reconstruct unknown hidden topology. They
specialize established linear observability, max-flow/min-cut, and spanning-tree
methods to a declared current-measurement contract. Conservation-based sensor
placement and measured cotrees have substantial prior literature, including
Bagajewicz and S\'anchez (2000) and He (2013); no priority claim is made for these
graph primitives. The purpose is to connect them explicitly to channel access,
sharp error guarantees, and physical realizability.

\section{The observation contract}

Let $G=(V,E)$ be a finite undirected multigraph without self-loops. Each physical
channel $e$ has an arbitrary reference orientation $u_e\to v_e$. Parallel edges
represent different transition mechanisms between the same pair of states.
The incidence matrix $B\in\mathbb R^{V\times E}$ has $-1$ at a channel's tail
and $+1$ at its head. A stationary current vector belongs to
\begin{equation}
 \mathcal J=\ker B.
 \label{eq:graph-current-class}
\end{equation}
Reference orientation is bookkeeping: every component of $j\in\mathcal J$
may be positive, negative, or zero. Initially, the model class imposes no
activity, rate, affinity, or current ceiling. Such additional restrictions can
change the sharp uncertainty bounds below.
The minimum-cut formulas require an incidence matrix. A general chemical
stoichiometric balance matrix need not have that form; it requires the more
general linear-identification and robust-recovery formulation rather than an
automatic application of these graph formulas.

A sensor for channel $s$ supplies either its exact native current $j_s$ or a
number $z_s$ with a justified bound $|z_s-j_s|\leq\epsilon_s$. The apparatus
need not observe the full state trajectory. This distinction matters: complete
state-resolved jump trajectories already identify aggregate transition rates
and their aggregate currents, subject to statistical uncertainty. The problem
here concerns sparse current access, including local intervention sensors that
report selected physical-channel currents.

Every impossibility or minimax statement is relative to the retained sensor
records. Richer raw observations may contain additional information. A sensor's
construction, calibration, and resource accounting must be supplied separately;
assigning an edge a cost does not by itself give that cost an energy meaning.

For a vertex set $U$, write
\begin{equation}
 \sigma_U(e)=\mathbf 1\{u_e\in U\}-\mathbf 1\{v_e\in U\}.
\end{equation}
Summing $Bj=0$ over the vertices of $U$ gives
\begin{equation}
 \sum_{e\in E}\sigma_U(e)j_e=0.
 \label{eq:graph-cut-conservation}
\end{equation}
Fix a target $e_0:u\to v$. For $u\in U$ and $v\notin U$, let
$C(U)=\{e\ne e_0:\sigma_U(e)\ne0\}$. This is the corresponding $u$--$v$
cut in $G\setminus\{e_0\}$.

\section{Exact identification and the minimum additional access}

\begin{theorem}[A target current requires a measured separating cut]
\label{thm:graph-target-identification}
Let $S\subseteq E\setminus\{e_0\}$ be the channels observed exactly. The
following conditions are equivalent:
\begin{enumerate}
 \item $j_S=j'_S$ implies $j_{e_0}=j'_{e_0}$ for every $j,j'\in\mathcal J$.
 \item $e_0$ is a bridge in $G\setminus S$.
 \item $S$ separates $u$ from $v$ in $G\setminus\{e_0\}$.
 \item Some $U$ containing $u$ and excluding $v$ satisfies $C(U)\subseteq S$.
\end{enumerate}
Whenever these conditions hold,
\begin{equation}
 j_{e_0}=-\sum_{s\in C(U)}\sigma_U(s)j_s.
 \label{eq:graph-target-estimator}
\end{equation}
\end{theorem}

\begin{proof}
Suppose a path from $v$ to $u$ remains after deleting $S$ and $e_0$. Send one
unit of signed flow along that path, and one unit along $e_0$ from $u$ to $v$.
The resulting vector $h$ obeys $Bh=0$, $h_S=0$, and $h_{e_0}=1$. Thus $j=0$
and $j'=h$ have identical sensor records and different target currents.
Identification is impossible.

If no such path remains, let $U$ be the component of $u$ after those deletions.
Then $v\notin U$, and every original crossing edge except $e_0$ belongs to
$S$. Equation~\eqref{eq:graph-cut-conservation}, with $\sigma_U(e_0)=1$,
gives~\eqref{eq:graph-target-estimator}. It determines the target from the
record. The existence of the alternative path is exactly the failure of $e_0$
to be a bridge in $G\setminus S$, proving all equivalences.
\end{proof}

An original bridge has stationary current zero. The theorem then requires no
sensor: conservation already fixes its current. Parallel channels must not be
collapsed before applying the theorem. A second unsensed channel between $u$
and $v$ supplies an alternative path of length one and hence an ambiguity.

\begin{corollary}[Minimum additional sensing cost]
\label{cor:graph-minimum-cost}
Give each eligible nontarget channel a nonnegative acquisition cost $c_s$.
Other channels cannot be selected. The minimum total cost of an exact indirect
target measurement is
\begin{equation}
 \min_{U:\,u\in U,\,v\notin U}\ \sum_{s\in C(U)}c_s,
 \label{eq:graph-installation-cut}
\end{equation}
where a cut containing an ineligible channel has infinite cost. If direct
measurement of $e_0$ is also allowed, take the smaller of its cost and the
indirect optimum.
\end{corollary}

\begin{proof}
Every successful sensor set contains $C(U)$ for some admissible $U$, by
Theorem~\ref{thm:graph-target-identification}. Nonnegative costs make its
total cost at least that of this cut. Conversely, measuring exactly the edges
of a minimum cut supplies the estimator~\eqref{eq:graph-target-estimator}.
\end{proof}

This is an ordinary minimum-cut computation with a scientific interpretation.
Its optimality is for the supplied sensor menu and additive cost model. It
does not assert that the menu contains every physically possible experiment.

\section{A sharp error bound, including nonlinear estimators}

Let $S$ now supply noisy current estimates
\begin{equation}
 z_s=j_s+\eta_s,\qquad |\eta_s|\leq\epsilon_s,\qquad s\in S.
 \label{eq:graph-box-record}
\end{equation}
The full Cartesian error box is allowed. This is an adversarial uncertainty
set, not an assumption of probabilistic independence. Define the deterministic
minimax radius
\begin{equation}
 R_S=\inf_f\ \sup_{\substack{j\in\ker B\\|\eta_s|\leq\epsilon_s}}
       \left|f(j_S+\eta)-j_{e_0}\right|,
 \label{eq:graph-minimax-definition}
\end{equation}
where $f$ ranges over all estimators, including nonlinear ones.

\begin{theorem}[The exact radius is an error-weighted cut]
\label{thm:graph-exact-radius}
Under~\eqref{eq:graph-current-class} and~\eqref{eq:graph-box-record},
\begin{equation}
 R_S=\min_{\substack{U:\,u\in U,\,v\notin U\\C(U)\subseteq S}}
                 \sum_{s\in C(U)}\epsilon_s.
 \label{eq:graph-exact-radius}
\end{equation}
The value is $+\infty$ when there is no admissible cut. For a finite minimum,
the signed cut-sum estimator attains the radius globally.
\end{theorem}

\begin{proof}
For any admissible cut, use
$f(z)=-\sum_{s\in C(U)}\sigma_U(s)z_s$. Subtracting the exact conservation
identity shows that its error is a signed sum of the $\eta_s$. Its absolute
value is at most $\sum_{s\in C(U)}\epsilon_s$. Minimizing over cuts gives the
upper bound.

For the matching lower bound, remove $e_0$. Give each sensor edge capacity
$\epsilon_s$ in either direction, and every unobserved edge infinite
capacity. If a finite separating cut exists, undirected max-flow/min-cut gives
a signed flow of value $R$ from $v$ to $u$, where $R$ is the right side
of~\eqref{eq:graph-exact-radius}. Close the flow with current $R$ on $e_0$.
The resulting circulation $h$ satisfies
\begin{equation}
 Bh=0,\qquad h_{e_0}=R,\qquad |h_s|\leq\epsilon_s\quad(s\in S).
 \label{eq:graph-dual-witness}
\end{equation}
Both $j=h$ with errors $\eta=-h_S$ and $j=-h$ with errors $\eta=h_S$ produce
the identical record $z=0$. Their target values are $R$ and $-R$. Any single
estimate $f(0)$ has error at least $R$ in one of the two cases. This lower
bound applies to every estimator and matches the upper bound.

For completeness, the infinite-capacity construction can be made finite. If
there is no entirely unobserved $v$--$u$ path, assign every unobserved edge a
capacity strictly larger than $\sum_{s\in S}\epsilon_s$. At least one cut
uses only sensor edges, so a minimum cut never crosses an unobserved edge.
The finite-capacity max-flow theorem therefore gives exactly the required
flow. If an entirely unobserved path does exist, the unit circulation in the
proof of Theorem~\ref{thm:graph-target-identification} can be scaled without
bound while leaving every sensor record unchanged. Then $R_S=+\infty$.
\end{proof}

The theorem explains why averaging redundant sensors need not improve a
worst-case box guarantee. It does not prohibit statistical improvement under a
specified noise law or an uncertainty set that excludes the least-favorable
error combinations. Known rate or activity bounds can also reduce the sharp
radius by restricting the allowed currents.

\paragraph{A certificate requiring no optimization.}
A finite certificate consists of a vertex cut $U$ and a circulation $h$. A
checker verifies $C(U)\subseteq S$, the identities and bounds
in~\eqref{eq:graph-dual-witness}, and
$h_{e_0}=\sum_{s\in C(U)}\epsilon_s$. These checks establish matching lower
and upper bounds without trusting or rerunning the search algorithm. An
infinite certificate instead supplies $Bh=0$, $h_{e_0}=1$, and $h_S=0$.

\section{Aggregate access and the internal channel obstruction}

Partition the physical channels into bundles $\mathcal B$, one for each
unordered pair of endpoint states. Choose one canonical orientation per bundle.
Let $A$ take signed bundle sums, so $k=Aj$ is the vector of aggregate currents
and
\begin{equation}
 B=\bar B A,
 \label{eq:graph-bundle-factorization}
\end{equation}
where $\bar B$ is the incidence matrix of the underlying simple graph.
Suppose bundle totals are known exactly for $T\subseteq\mathcal B$, and
individual channels are sensed exactly for $S\subseteq E$. These aggregate
readings are treated as already available. If their calibration errors are
nonzero, a robust linear-observation formulation is needed instead of silently
using the exact identities below.

Write $r_b=|b\setminus S|$ for the number of unsensed channels in bundle $b$.
Form a graph $H$ on all vertices $V$, including isolated vertices, with one
edge for each bundle $b\notin T$ having $r_b>0$. Let $c(H)$ denote its number
of connected components.

\begin{theorem}[Two distinct sources of current ambiguity]
\label{thm:graph-mixed-access}
The homogeneous ambiguity space
\begin{equation}
 \mathcal N=\{h:Bh=0,\ A_T h=0,\ h_S=0\}
\end{equation}
has dimension
\begin{equation}
 \dim\mathcal N
 =\sum_{b\in\mathcal B}\max(r_b-1,0)
       +|E(H)|-|V|+c(H).
 \label{eq:graph-mixed-dimension}
\end{equation}
An unsensed target $e_0$ in bundle $b_0$ is identifiable if and only if
\begin{equation}
 r_{b_0}=1
 \quad\hbox{and}\quad
 \bigl(b_0\in T\ \hbox{or $b_0$ is a bridge of $H$}\bigr).
 \label{eq:graph-mixed-target}
\end{equation}
\end{theorem}

\begin{proof}
After imposing $h_S=0$, each nonempty unsensed bundle has $r_b$ free channel
coordinates. Taking its signed sum maps these coordinates onto one aggregate
coordinate. The kernel of that map has dimension $r_b-1$. Because bundles have
disjoint channel coordinates, their internal kernels form a direct sum of
dimension $\sum_b\max(r_b-1,0)$.

Known aggregate variations vanish. All remaining aggregate variations must
satisfy conservation on $H$, by~\eqref{eq:graph-bundle-factorization}; hence
they belong to $\ker B_H$. Conversely, choose one unsensed representative
channel in each nonempty bundle. Any circulation on $H$ can be lifted by
assigning its aggregate current, with the appropriate orientation sign, to
that representative and assigning zero to every other channel. Known bundle
variations remain zero, and $Bh=\bar BAh=0$. The map from $\mathcal N$ onto
$\ker B_H$ is therefore surjective, with exactly the internal kernel described
above. The dimension formula follows from
$\dim\ker B_H=|E(H)|-|V|+c(H)$.

If $r_{b_0}\geq2$, one can change the target and compensate on another
unsensed channel in the same bundle. Choose the signs so that their aggregate
variation is zero. This changes $h_{e_0}$ while preserving all supplied
measurements and conservation, so the target is unidentifiable. If
$r_{b_0}=1$, the target variation equals the aggregate variation of $b_0$ up
to orientation. It vanishes when $b_0\in T$. Otherwise, it vanishes for every
circulation on $H$ precisely when $b_0$ is a bridge, by
Theorem~\ref{thm:graph-target-identification}. This proves both directions
of~\eqref{eq:graph-mixed-target}.
\end{proof}

The first term in~\eqref{eq:graph-mixed-dimension} is ambiguity between
physical channels sharing the same endpoints. The second is circulation
ambiguity among unknown aggregate edges. The decomposition need not select a
preferred representative channel, but its dimension and target criterion are
independent of that choice.

In particular, if every aggregate current is known, then $\ker A\subseteq\ker
B$ by~\eqref{eq:graph-bundle-factorization}. Conservation supplies no further
constraint on the internal split. It cannot distinguish two physical channels
whose currents change oppositely while leaving the aggregate unchanged. This
is a necessary-access obstruction: some measurement or intervention must
distinguish those channels, unless additional physical assumptions already
fix their relation.

\section{Constructive completion with mixed access}

\begin{proposition}[Minimum indirect access for one target]
\label{prop:graph-mixed-completion}
Assume all nontarget channels can be sensed at nonnegative costs $c_e$, and
the exact aggregate readings in $T$ are already available. A minimum-cost
indirect sensor set for target $e_0\in b_0$ is obtained as follows:
\begin{enumerate}
 \item Sense every other channel in $b_0$.
 \item If $b_0\in T$, stop.
 \item Otherwise, build the simple graph of bundles not in $T$, remove
       $b_0$, and give each other bundle the cost
       $C_b=\sum_{e\in b}c_e$.
 \item Find a minimum cut between the endpoints of $b_0$ and sense every
       physical channel in every bundle crossing that cut.
\end{enumerate}
\end{proposition}

\begin{proof}
Every other channel in $b_0$ must be sensed: otherwise the internal two-channel
ambiguity in Theorem~\ref{thm:graph-mixed-access} survives. These costs are
unavoidable. If the target aggregate is known, the forced sensors already
determine the target by subtraction.

If the target aggregate is unknown, it must become a bridge of the residual
unknown-bundle graph. A nontarget bundle disappears from that graph exactly
when all its physical channels are sensed, at cost $C_b$. Partially sensing a
bundle leaves at least one channel able to carry its arbitrary aggregate
variation. By Theorem~\ref{thm:graph-target-identification}, every successful
design must therefore remove the bundles of a separating cut. Its additional
cost is at least the minimum cut cost. Sensing all channels on that minimum cut
makes the target a bridge and attains the bound. These bundles are disjoint
from $b_0$, so their costs add to the forced parallel-channel costs.
\end{proof}

\begin{proposition}[Minimum access for every physical-channel current]
\label{prop:graph-all-targets}
Suppose every channel has a nonnegative sensing cost. Define
$w_b=\max_{e\in b}c_e$. The minimum total cost for recovering every current,
given exact aggregate access on $T$, is
\begin{equation}
 \sum_{e\in E}c_e
 -\sum_{b\in T}w_b
 -\max_{F\text{ forest in }\mathcal B\setminus T}\ \sum_{b\in F}w_b.
 \label{eq:graph-all-target-optimum}
\end{equation}
A maximum-weight spanning forest of the unknown-bundle graph supplies an
optimal design.
\end{proposition}

\begin{proof}
Full recovery requires zero ambiguity dimension in
Theorem~\ref{thm:graph-mixed-access}. Thus at most one channel may remain
unsensed in each bundle, and the unknown bundles retaining such a channel must
form a forest. Conversely, these two conditions make both terms of the
dimension formula zero and hence suffice.

To maximize saved sensing cost, retain a maximum-cost unsensed channel in
each known-total bundle; these choices impose no cycle constraint. For an
unknown-total bundle that is retained, the best saved cost is also $w_b$.
The set of retained unknown bundles must be a forest, so a maximum-weight
forest maximizes their total savings. Subtracting both savings from the cost
of sensing all channels gives~\eqref{eq:graph-all-target-optimum}. Nonnegative
weights allow a maximizing forest to be extended to a spanning forest in
each component without decreasing its weight.
\end{proof}

With no aggregate access, this reduces to the familiar measured-cotree rule
on the physical multigraph. With every aggregate known, one maximum-cost
channel per bundle can be left unsensed, but all other channel splits require
individual information. The algorithms are exact for their declared access
and price assumptions; uncertain aggregate measurements or unavailable
channels require the corresponding constraints to be retained explicitly.

\section{The ambiguity witnesses are realizable by positive rates}

Allowing signed currents does not require negative transition rates.

\begin{proposition}[Realization at a prescribed stationary distribution]
\label{prop:graph-physical-realization}
Let $j\in\ker B$ be finite, and let $\pi$ be any strictly positive probability
distribution on $V$. Choose $t_e>|j_e|/2$ for every channel and define
\begin{equation}
 q_e(u_e,v_e)=\frac{t_e+j_e/2}{\pi_{u_e}},\qquad
 q_e(v_e,u_e)=\frac{t_e-j_e/2}{\pi_{v_e}}.
 \label{eq:graph-rate-realization}
\end{equation}
The resulting bidirectional channel rates are strictly positive. Their sum
defines a continuous-time Markov chain with stationary distribution $\pi$ and
channel currents exactly $j$. If the support graph is connected, the chain is
irreducible.
\end{proposition}

\begin{proof}
The strict traffic inequality makes both numerators positive. The stationary
forward and reverse channel fluxes are respectively $t_e+j_e/2$ and
$t_e-j_e/2$, whose difference is $j_e$. Summing their net contributions at
each vertex gives $Bj=0$, so incoming and outgoing stationary fluxes balance.
Setting generator diagonals to minus outgoing rates therefore gives
$\pi Q=0$. Connected bidirectional support makes every state reachable from
every other state and hence makes the chain irreducible.
\end{proof}

A finite pair of least-favorable currents $h$ and $-h$ can use the same $\pi$
and the same traffic $t_e>|h_e|/2$. Thus the sharp-error witnesses can have
identical stationary occupations and positive finite rates. An unbounded
impossibility result uses arbitrarily large finite rates across the model
class, not an infinite rate in an individual model. A fixed physical device
may impose traffic, fuel, or affinity limits that restrict this class. Fixed
chemical stoichiometry and fixed reservoirs also introduce constraints not
encoded by conservation alone.

\section{Worked example and reproducible interface}

Consider states $A,B,C,D$ and six physical channels:
\begin{equation}
 e_0:A\to B,\quad p:A\to B,\quad
 a:A\to C,\quad b:B\to C,\quad c:C\to D,\quad d:D\to A.
\end{equation}
Suppose only $p,a,d$ are measured, with respective error radii
$1/100$, $1/50$, and $3/100$. The cut $U=\{A\}$ gives
\begin{equation}
 \widehat j_{e_0}=-z_p-z_a+z_d,\qquad R=\frac{3}{50}.
\end{equation}
The following matching circulation certifies sharpness:
\begin{equation}
 (h_{e_0},h_p,h_a,h_b,h_c,h_d)
 =\left(\frac3{50},-\frac1{100},-\frac1{50},
                  \frac1{20},\frac3{100},\frac3{100}\right).
\end{equation}
It obeys conservation, reaches the target error radius, and respects all
sensor error bounds. Its two signs can be realized by
\eqref{eq:graph-rate-realization} with common stationary occupations.

If all five bundle totals are instead known exactly, one ambiguity remains:
the split between $e_0$ and $p$. Reading $p$ resolves it. Conservation alone
cannot. For a separate exact sensing-cost example, assign costs
$c_p=2,c_a=3,c_b=4,c_c=5,c_d=6$. The minimum indirect design for $e_0$ measures
$p$ and $b$, at total cost $6$. Conservation gives $j_{e_0}=-j_p+j_b$.
This acquisition-price example and the preceding error-radius example use
different declared inputs; their numbers must not be interpreted as the same
resource quantity.

The implementation is in \texttt{code/graph\_access.py}. It accepts exact
integers, rational strings, or \texttt{Fraction} values and rejects floating
point capacities. Its main interfaces are:
\begin{description}
 \item[\texttt{target\_cut}] Finds a minimum cut for supplied acquisition
 costs or sensor-error capacities, returning a cut and matching circulation.
 \item[\texttt{verify\_target\_cut}] Checks the certificate directly,
 without invoking an optimization algorithm.
 \item[\texttt{mixed\_observability}] Returns the internal-channel and
 network-cycle ambiguity dimensions and all identifiable target channels.
 \item[\texttt{target\_completion}] Returns a minimum indirect sensor set
 with exact aggregate access; direct measurement of the target is excluded.
 \item[\texttt{all\_target\_completion}] Returns an optimal full-current
 design using the maximum-weight forest construction.
 \item[\texttt{realize\_currents}] Constructs positive rational channel
 rates at a prescribed strictly positive stationary distribution.
\end{description}
The rational example certificates and rate realization are saved in
\texttt{results/graph\_access.json}. Reproduction commands are
\begin{verbatim}
python code/graph_access.py
python -m unittest discover -s tests -p test_graph_access.py -v
\end{verbatim}

Author tests compare cut values with exhaustive vertex partitions, mixed
observability with an independent rational rank calculation, and both sensor
designs with exhaustive subsets. Separate internal review additionally checked
2,197 mixed-observation configurations, 65 completion instances, and 160 cut
cases. These checks support the implementation; the proofs establish the
general results. External peer review and experimental validation remain
separate requirements.
